% ============================================================================
% kapitel/einleitung.tex
% ============================================================================

\section{Das Problem der modernen Physik}

Die moderne Physik steht vor einem Rätsel. Sie verfügt mit der Quantenmechanik und der Allgemeinen Relativitätstheorie über zwei der erfolgreichsten Theorien der Wissenschaftsgeschichte. Aber sie sind fundamental inkompatibel.

Die vorgeschlagenen Lösungen – Stringtheorie, Schleifenquantengravitation, dunkle Materie, dunkle Energie – haben das Problem nicht gelöst.

Es fehlt ein einheitliches Fundament.

\section{Die Lösung: IWT und Q}

Dieses Dokument stellt die Informations-Weber-Theorie (IWT) vor.

Sie hat nur eine fundamentale Größe: \textbf{Information} \( I \).

Aus der Information emergiert alles andere: Raum, Zeit, Materie, Energie, Gravitation, Quantenmechanik, Elektrodynamik.

Die organisierende Kraft ist \textbf{Q}: die Quelle der Selbstorganisation.

\section{Zwei Quellen}

Dieses Dokument ruht auf zwei Säulen:

\begin{enumerate}
    \item \textbf{Die Theorie:} Axiome, Herleitungen, mathematische Struktur
    \item \textbf{Die Simulation:} Numerische Bestätigung, Ergebnisse, Visualisierung
\end{enumerate}

Die Theorie sagt, was sein muss. Die Simulation zeigt, dass es ist.

\section{Die zentrale Aussage}

\textbf{Q ist die Quelle.}

Alles emergiert aus Q.

\begin{itemize}
    \item Die DBT emergiert aus Q (Bohm-Potential)
    \item Die WG emergiert aus Q (Gravitation)
    \item Die WED emergiert aus Q (Elektrodynamik)
\end{itemize}

Q ist das Prinzip der Selbstorganisation.

Q ist die Quelle der Physik.